\begin{exo}
	Une urne contient deux boules vertes ($V_1$ et $V_2$) et trois boules noires ($N_1$, $N_2$ et $N_3$) toutes indiscernables au toucher.

	On tire successivement et \textbf{sans remise} deux boules dans l'urne et on considère les événements :
	\begin{itemize}[label=$\bullet$]
		\item $A$ : \og{} les deux boules sont de même couleur \fg{} ;
		\item $B$ : \og{} au moins une boule est verte \fg{}.
	\end{itemize}

	Écrire chaque événement sous forme d'un ensemble des issues possibles et calculer leur probabilité.
\end{exo}
